%% kapa: exam=1 round=19 year=2008 period=1 form=A subject=civil_law subject_ko=민법 seq=38 number=38
%% source: Q-Net 감정평가사 기출문제 (cst00309) artlSeq=1001962 file=extracted/1008/1001962/2008년 제19회 감정평가사 1차시험 기출문제/1교시.hwp
%% conversion: kordoc 4.24.51 -> tools/gen1.py (auto); review=proofread-sonnet; flags=none
\begin{question}{38}
양도담보에 관한 설명으로 틀린 것은? (다툼이 있는 경우에는 판례에 의함)
\begin{choices}
  \item 양도담보가 설정된 경우 특별한 사정이 없는 한 목적부동산에 대한 사용․수익권은 양도담보설정자에게 있으므로, 목적부동산을 임대할 권한은 양도담보설정자에게 있다.
  \item 사용․수익할 수 있는 정당한 권한이 있는 양도담보설정자로부터 그 권한을 승계한 자에 대하여 양도담보권자는 사용․수익을 하지 못한 것을 이유로 임료상당의 손해배상이나 부당이득반환을 청구할 수 없다.
  \item 부동산의 양도담보권설정자는 그 부동산의 등기명의가 양도담보권자 앞으로 되어 있더라도 그 부동산의 불법점유자인 제3자에 대하여는 실질적 소유자임을 주장하여 불법점유의 상태의 배제권을 행사할 수 있다.
  \item 양도담보권자는 담보권의 실행을 위하여 담보채무자가 아닌 제3자에 대하여도 담보물의 인도를 구할 수 있다.
  \item 점유개정의 방법으로 동산에 대하여 이중으로 양도담보설정계약이 체결된 경우, 뒤에 설정계약을 체결한 자가 양도담보권을 선의취득한다.
\end{choices}
\end{question}
\answer{5}
